Priyansh wrote the trace-generation script and the Python reference model used to validate the cache. The generator, \texttt{scripts/gen\_traces.py}, is seeded and produces five 10{,}000-access traces in the \texttt{R~0xADDR} / \texttt{W~0xADDR~0xDATA} format: sequential word reads, uniform random over 1\,MB, a loop over a 512\,B working set, a 64\,B stride, and a matrix-multiply-like pattern. The reference model, \texttt{scripts/cache\_model.py}, implements the frozen specification (1\,KB, 16\,B blocks, 64 lines, tag $[31:10]$, index $[9:4]$, offset $[3:0]$, write-through with no write allocate) and reports accesses, hits, misses, reads, writes, and hit rate per trace; it also accepts cache and block sizes as options for the Phase~II sweeps. Its hit rates (75.00\% sequential, 0.12\% random, 99.68\% loop, 0.00\% stride, 88.85\% matmul) agree with the expected behavior of each pattern and serve as the golden values for the RTL counter cross-check. In Phase~II, Priyansh is responsible for the cross-check, the hit-rate and AMAT tables, and the size and block-size sweep plots.
